\documentclass[11pt,a4paper]{article}
\usepackage[margin=1in]{geometry}
\usepackage{kotex}
\usepackage{graphicx}
\usepackage{booktabs}
\usepackage{hyperref}
\title{Attention Is All You Need 간략 보고서}
\author{OpenClaw Auto Brief}
\date{\today}
\begin{document}
\maketitle

\section{한 줄 요약}
\textit{Attention Is All You Need} (Vaswani et al., 2017)은 RNN/CNN 없이 \textbf{Self-Attention}만으로 시퀀스 변환을 수행하는 Transformer를 제안했고,
이는 현대 LLM의 핵심 아키텍처 표준이 되었다.

\section{문제의식}
기존 Seq2Seq(RNN 기반)는 다음 한계가 있었다.
\begin{itemize}
  \item 장거리 의존성 학습이 어렵다.
  \item 순차 처리 특성상 학습 병렬화가 제한적이다.
  \item 학습 시간/비용이 크다.
\end{itemize}
논문은 이를 해결하기 위해 토큰 간 관계를 직접 계산하는 Attention 중심 구조를 제시했다.

\section{핵심 아이디어}
\subsection{Transformer 인코더-디코더}
모델은 인코더/디코더 블록을 쌓는 구조이며, 각 블록은 다음으로 구성된다.
\begin{itemize}
  \item Multi-Head Self-Attention
  \item Position-wise Feed-Forward Network
  \item Residual Connection + Layer Normalization
\end{itemize}

\subsection{Scaled Dot-Product Attention}
\textbf{Q, K, V}를 사용해 토큰 간 중요도를 계산한다.
\[
\mathrm{Attention}(Q,K,V)=\mathrm{softmax}\left(\frac{QK^T}{\sqrt{d_k}}\right)V
\]
$\sqrt{d_k}$로 스케일링해 큰 차원에서 softmax가 과도하게 뾰족해지는 문제를 완화한다.

\subsection{Multi-Head Attention}
하나의 attention 대신 여러 head를 병렬로 사용해 서로 다른 관계(문법적/의미적/장거리 의존성)를 동시에 포착한다.

\subsection{Positional Encoding}
RNN이 없으므로 순서 정보를 명시적으로 주입한다. 사인/코사인 기반 positional encoding을 더해 토큰 순서를 표현한다.

\section{주요 결과와 의의}
WMT 2014 번역 과제(EN-DE, EN-FR)에서 강력한 성능을 보였고,
동시에 학습 병렬화 효율을 크게 개선했다. 핵심 의의는 다음과 같다.
\begin{itemize}
  \item 시퀀스 모델링의 중심을 recurrence에서 attention으로 이동
  \item 대규모 사전학습(Pretraining) 기반 언어모델 발전의 토대 제공
  \item BERT, GPT 계열을 포함한 현대 LLM 아키텍처의 직접적 기반
\end{itemize}

\section{제한점과 후속 연구}
원 논문 이후 알려진 대표적 한계:
\begin{itemize}
  \item Self-Attention의 $O(n^2)$ 메모리/연산 비용
  \item 긴 문맥에서의 비용 증가
\end{itemize}
이를 해결하기 위해 Sparse/Linear Attention, FlashAttention, 긴 컨텍스트 최적화 연구가 활발히 진행됐다.

\section{실무 관점 정리}
\begin{enumerate}
  \item Transformer 이해는 LLM 이해의 출발점이다.
  \item 모델 성능뿐 아니라 \textbf{학습/추론 효율}이 중요하다.
  \item Attention 메커니즘과 토크나이징/포지셔널 설계는 실제 성능에 큰 영향을 준다.
\end{enumerate}

\section*{참고 문헌}
Vaswani, A., et al. (2017). \textit{Attention Is All You Need}. NeurIPS 2017.

\end{document}
